\section{Policy Comparison}\label{sec:policies}

With the raw traces characterized, we now compare six policies --- LRU, FIFO,
CLOCK, S3-FIFO, SIEVE, and W-TinyLFU --- head-to-head. Unless otherwise
noted, every curve, number, and significance claim in this section rests on
$5$-seed replay-Zipf sweeps across $\alpha \in \{0.6, 0.7, \ldots, 1.2\}$ and
cache fractions $\{0.001, 0.005, 0.010, 0.020, 0.050, 0.100\}$ on both
workloads, with Welch's two-sample $t$-test (\texttt{scipy.stats.ttest\_ind},
\texttt{equal\_var=False}) against LRU as the reference policy at a $p < 0.05$
threshold. The Zipf $\alpha$ estimates themselves are MLE-fit per
\cite{clauset2009powerlaw}.

\subsection*{Headline result}

\textbf{W-TinyLFU wins 11 of 12 miss-ratio cells} ($6$ cache fractions
$\times$ $2$ workloads). The one exception is Congress at
$\mathrm{cache\_frac}=0.020$ where SIEVE edges W-TinyLFU by less than $1\sigma$.
Figure~\ref{fig:opening} (page~1) shows both workloads side-by-side with
$\pm 1\sigma$ bands across $5$ seeds. Note the visual asymmetry: Congress
curves all collapse into a narrow band at high $\alpha$ while Court curves
fan out into clean strata --- the same policies, the same sweep grid, yet
different separability.

\subsection*{Why W-TinyLFU dominates}

W-TinyLFU is the only policy in our set with \emph{explicit frequency
tracking}. Three mechanisms account for the dominance:

\begin{itemize}[leftmargin=*,itemsep=0pt,topsep=2pt]
  \item A $4$-bit Count-Min Sketch tracks approximate per-key access
    frequency across the whole cache.
  \item An admission filter requires new candidates to out-frequency the
    current victim before they are allowed in.
  \item A $1\%$ window-LRU gives one-hit wonders a chance to prove
    recurrence rather than immediately contaminating the main region.
\end{itemize}

Every other policy we tested uses only \emph{recency} signals. Under Zipf
popularity with moderate-to-high $\alpha$, a key's long-run frequency is a
better predictor of its next-access probability than its most-recent-access
time --- and only W-TinyLFU's CMS tracks that directly.

\textbf{Where the win is largest:} Small Cache (cache\_frac$=0.001$) on
Court --- W-TinyLFU beats LRU by $12.9$\,pp ($0.831$ vs.\ $0.960$). When the
cache holds $\sim 10$ objects out of a $15{,}018$-unique-key workload,
every admission is life-or-death: one wrong admission of a one-hit wonder
forever evicts a repeat-access key. TinyLFU's admission filter is literally
the thing preventing this disaster. \textbf{Where the win is smallest:}
large caches on Congress (cache\_frac$=0.100$) where roughly $1{,}900$
objects fit and the cache is big enough to absorb most of the hot set
regardless of eviction policy --- there the W-TinyLFU-vs-LRU gap shrinks
to about $5$\,pp. Caffeine's production configuration
(\cite{caffeine2024}) motivates the specific $1\%$ window / $99\%$ main
split we use.

\subsection*{Per-workload alpha sensitivity}

\begin{figure}[h]
\centering
\begin{minipage}{0.48\textwidth}
\includegraphics[width=\textwidth]{../results/congress/figures/alpha_sensitivity.pdf}
\end{minipage}\hfill
\begin{minipage}{0.48\textwidth}
\includegraphics[width=\textwidth]{../results/court/figures/alpha_sensitivity.pdf}
\end{minipage}
\caption{Alpha sensitivity at $\mathrm{cache\_frac}=0.01$ --- Congress
(left) vs.\ Court (right). W-TinyLFU dominates LRU monotonically across
$\alpha \in \{0.6,\ldots,1.2\}$ on both workloads, but the margin on Court
is larger and less dependent on $\alpha$ than on Congress.}
\label{fig:alpha_sens}
\end{figure}

Figure~\ref{fig:alpha_sens} plots miss-ratio against $\alpha$ at
$\mathrm{cache\_frac}=0.01$ for every policy, separately per workload.
\textbf{W-TinyLFU beats LRU monotonically by $7.84\%$--$21.55\%$ across
$\alpha \in \{0.6, \ldots, 1.2\}$ on Congress}, with the advantage growing
as skew rises (this matches TinyLFU theory on OHW filtering). The
Congress-side Welch's $t$-test delivers $p \in [1.2\times 10^{-6},\;
8.7\times 10^{-8}]$ at every $\alpha$ in the sweep --- dominance is
statistically significant at a $p < 10^{-6}$ threshold at every point, not
a small-effect-size artifact. Court shows larger absolute effect sizes
but also materially higher seed-to-seed variance (Court $\sigma \approx
0.045$ vs.\ Congress $\sigma \approx 0.007$ at W-TinyLFU,
$\mathrm{cache\_frac}=0.01$); the Court W-TinyLFU-vs-LRU $p$-value at
$\alpha = 1.2$ is $0.06$, non-significant despite an $8$-pp effect, which
motivates our insistence on multi-seed aggregation for heavy-tailed
workloads.

\subsection*{Policy ordering is robust across workloads}

Both workloads produce nearly identical policy orderings at
$\mathrm{cache\_frac}=0.01$:

\begin{itemize}[leftmargin=*,itemsep=0pt,topsep=2pt]
  \item \textbf{Congress:} W-TinyLFU $\approx$ SIEVE $>$ S3-FIFO $>$ CLOCK $>$ LRU $>$ FIFO
  \item \textbf{Court:} W-TinyLFU $>$ SIEVE $>$ S3-FIFO $>$ CLOCK $>$ LRU $>$ FIFO
\end{itemize}

The \emph{only} pairwise swap across the two workloads is W-TinyLFU
$\leftrightarrow$ SIEVE: statistically tied on Congress (gaps always within
$\pm 0.005$, less than $1\sigma$) and cleanly separated by $4$--$5$\,pp on
Court. Every other pair preserves its ordering. This is a finding, not a
non-finding: policy \emph{choice} is robust across a change in workload, but
the \emph{magnitudes} of the gaps are not --- they track workload
characteristics (raw $\alpha$, size-distribution tail) in ways the policy
mechanisms themselves do not fully anticipate. The mechanism behind the
SIEVE/W-TinyLFU swap is the subject of \S\ref{sec:mechanism}.
